\subsubsection*{LDPolypVideo}

The LDPolypVideo dataset is a large-scale benchmark specifically curated for evaluating polyp detection in challenging real-world colonoscopy videos. It contains 40,266 annotated frames extracted from 160 clinical video sequences, where 33,884 frames include at least one polyp. In addition, over 860,000 unlabeled frames are provided for semi-supervised or unsupervised learning purposes.

To handle the complexity and redundancy of video data, a hybrid annotation method was employed: an expert first labeled the initial frame, after which a tracking algorithm (Kernelized Correlation Filters) predicted subsequent positions. These predictions were iteratively validated and corrected by endoscopy experts to ensure high annotation quality.

The dataset includes a range of difficult visual conditions such as motion blur, occlusions, specular reflections, and multiple polyps per frame, making it a robust benchmark for assessing algorithm generalizability. Four object detection models—YOLOv3, RetinaNet, Faster R-CNN, and CenterNet—were evaluated, all pre-trained on ImageNet and fine-tuned on a simpler dataset (CVC-ClinicDB) before being tested on LDPolypVideo.

Standard object detection metrics were reported, including precision, recall, F1-score, and F2-score. Results revealed a notable performance drop when models trained on simpler datasets were evaluated on LDPolypVideo, highlighting its complexity and real-world relevance. This dataset serves as a valuable testbed for assessing the robustness of real-time polyp detection systems.


Test Test